% Options for packages loaded elsewhere
\PassOptionsToPackage{unicode}{hyperref}
\PassOptionsToPackage{hyphens}{url}
%
\documentclass[
  10pt,
  letterpaper,
]{article}
\usepackage{amsmath,amssymb}
\usepackage{iftex}
\ifPDFTeX
  \usepackage[T1]{fontenc}
  \usepackage[utf8]{inputenc}
  \usepackage{textcomp} % provide euro and other symbols
\else % if luatex or xetex
  \usepackage{unicode-math} % this also loads fontspec
  \defaultfontfeatures{Scale=MatchLowercase}
  \defaultfontfeatures[\rmfamily]{Ligatures=TeX,Scale=1}
\fi
\usepackage{lmodern}
\ifPDFTeX\else
  % xetex/luatex font selection
\fi
% Use upquote if available, for straight quotes in verbatim environments
\IfFileExists{upquote.sty}{\usepackage{upquote}}{}
\IfFileExists{microtype.sty}{% use microtype if available
  \usepackage[]{microtype}
  \UseMicrotypeSet[protrusion]{basicmath} % disable protrusion for tt fonts
}{}
\makeatletter
\@ifundefined{KOMAClassName}{% if non-KOMA class
  \IfFileExists{parskip.sty}{%
    \usepackage{parskip}
  }{% else
    \setlength{\parindent}{0pt}
    \setlength{\parskip}{6pt plus 2pt minus 1pt}}
}{% if KOMA class
  \KOMAoptions{parskip=half}}
\makeatother
\usepackage{xcolor}
\usepackage[margin=0.85in]{geometry}
\usepackage{longtable,booktabs,array}
\usepackage{calc} % for calculating minipage widths
% Correct order of tables after \paragraph or \subparagraph
\usepackage{etoolbox}
\makeatletter
\patchcmd\longtable{\par}{\if@noskipsec\mbox{}\fi\par}{}{}
\makeatother
% Allow footnotes in longtable head/foot
\IfFileExists{footnotehyper.sty}{\usepackage{footnotehyper}}{\usepackage{footnote}}
\makesavenoteenv{longtable}
\usepackage{graphicx}
\makeatletter
\def\maxwidth{\ifdim\Gin@nat@width>\linewidth\linewidth\else\Gin@nat@width\fi}
\def\maxheight{\ifdim\Gin@nat@height>\textheight\textheight\else\Gin@nat@height\fi}
\makeatother
% Scale images if necessary, so that they will not overflow the page
% margins by default, and it is still possible to overwrite the defaults
% using explicit options in \includegraphics[width, height, ...]{}
\setkeys{Gin}{width=\maxwidth,height=\maxheight,keepaspectratio}
% Set default figure placement to htbp
\makeatletter
\def\fps@figure{htbp}
\makeatother
\setlength{\emergencystretch}{3em} % prevent overfull lines
\providecommand{\tightlist}{%
  \setlength{\itemsep}{0pt}\setlength{\parskip}{0pt}}
\setcounter{secnumdepth}{-\maxdimen} % remove section numbering
% definitions for citeproc citations
\NewDocumentCommand\citeproctext{}{}
\NewDocumentCommand\citeproc{mm}{%
  \begingroup\def\citeproctext{#2}\cite{#1}\endgroup}
\makeatletter
 % allow citations to break across lines
 \let\@cite@ofmt\@firstofone
 % avoid brackets around text for \cite:
 \def\@biblabel#1{}
 \def\@cite#1#2{{#1\if@tempswa , #2\fi}}
\makeatother
\newlength{\cslhangindent}
\setlength{\cslhangindent}{1.5em}
\newlength{\csllabelwidth}
\setlength{\csllabelwidth}{3em}
\newenvironment{CSLReferences}[2] % #1 hanging-indent, #2 entry-spacing
 {\begin{list}{}{%
  \setlength{\itemindent}{0pt}
  \setlength{\leftmargin}{0pt}
  \setlength{\parsep}{0pt}
  % turn on hanging indent if param 1 is 1
  \ifodd #1
   \setlength{\leftmargin}{\cslhangindent}
   \setlength{\itemindent}{-1\cslhangindent}
  \fi
  % set entry spacing
  \setlength{\itemsep}{#2\baselineskip}}}
 {\end{list}}
\usepackage{calc}
\newcommand{\CSLBlock}[1]{\hfill\break\parbox[t]{\linewidth}{\strut\ignorespaces#1\strut}}
\newcommand{\CSLLeftMargin}[1]{\parbox[t]{\csllabelwidth}{\strut#1\strut}}
\newcommand{\CSLRightInline}[1]{\parbox[t]{\linewidth - \csllabelwidth}{\strut#1\strut}}
\newcommand{\CSLIndent}[1]{\hspace{\cslhangindent}#1}
\ifLuaTeX
  \usepackage{selnolig}  % disable illegal ligatures
\fi
\usepackage{bookmark}
\IfFileExists{xurl.sty}{\usepackage{xurl}}{} % add URL line breaks if available
\urlstyle{same}
\hypersetup{
  hidelinks,
  pdfcreator={LaTeX via pandoc}}

\author{}
\date{}

\begin{document}

\section{PersonaMetrica: Stress-Testing Evaluation Protocols for
Long-Horizon Personal
Agents}\label{personametrica-stress-testing-evaluation-protocols-for-long-horizon-personal-agents}

\textbf{Anonymous submission draft --- NeurIPS Evaluations \& Datasets
framing}

\subsection{Abstract}\label{abstract}

Long-horizon personal assistants must do more than retrieve a previously
stated preference: they must determine which evidence still applies,
represent uncertainty when evidence conflicts, respect temporary scope,
and decide when a current belief is reliable enough to act on. Existing
memory and personalization evaluations often emphasize recall,
preference following, or state accuracy. We study a narrower evaluation
question: \textbf{which conclusions are stable once confidence is
allowed to control action, and which depend on the calibration protocol
itself?} We introduce \textbf{PersonaMetrica-Bench}, a controlled
evaluation framework that jointly measures current-state accuracy,
confidence calibration, action coverage, selective accuracy,
risk--coverage, and downstream decision utility under evolving
preferences. We also introduce \textbf{Personal Belief State (PBS)}, a
transparent evidence-weighted temporal reference system. The standard
release contains \textbf{500 synthetic users, 27,133 observations, and
13,500 decision queries} across ten preference domains. On the standard
controlled run PBS obtains \textbf{0.932 vs.~0.906 state accuracy},
\textbf{0.060 vs.~0.080 Brier}, and lower AURC (\textbf{0.0107
vs.~0.0208}) than a strong \texttt{time-aware-latest} baseline. However,
our most important result is negative: \textbf{fixed-threshold decision
utility is not protocol-invariant}. With native confidence, PBS appears
substantially better than time-aware-latest on 800-turn histories; after
every tracker receives the same low-capacity post-hoc calibration on
development seeds only, the long-horizon utility ranking reverses
(\textbf{0.588 PBS vs.~0.687 time-aware-latest}) while PBS retains
higher selective accuracy and slightly lower Brier. A separate frozen
audit over 75 protocol cells and all six released baselines finds that
time-aware-latest wins 45 cells, PBS 27, and first-mention 3; two
protocol cells choose different winners with probability \textbf{0.516},
mean complete-ranking Kendall \(\tau\) is \textbf{0.699}, and the most
discordant pair falls to \textbf{0.20}. A paired seed-bootstrap
diagnostic gives winner-change \textbf{0.516 {[}0.507, 0.538{]}}, and 13
leave-one-level-out grid perturbations all retain multiple winners. Thus
neither raw state accuracy nor fixed-threshold utility from unmatched
confidence scales is a safe standalone model-selection criterion. We
recommend reporting state accuracy, calibration, risk--coverage/AURC,
and utility only under an explicitly shared calibration-and-action
protocol. PersonaMetrica-Bench releases six baselines, development/test
splits, calibration and sensitivity analyses, component ablations, a
human-annotation interface, an external-model evaluation contract,
machine-readable metadata, and full reproduction artifacts. We
explicitly do not claim frontier-model or human-population
generalization from the controlled study.

\subsection{1. Introduction}\label{introduction}

Personalization is commonly framed as a memory problem: a user states a
preference, an assistant stores or retrieves it, and evaluation asks
whether the assistant later recalls or follows it. That abstraction
becomes incomplete once preferences evolve, evidence is noisy or
implicit, temporary exceptions occur, and memory is used to trigger
recommendations, reminders, tool actions, or other interventions.

Consider a user who generally prefers morning workouts, temporarily
switches to evenings during a two-week schedule disruption, and later
returns to mornings. A system that always uses the latest apparently
trustworthy fact can be highly accurate on aggregate state queries yet
become overly conservative or overconfident in exactly the cases where
action requires calibrated uncertainty. A second system can maintain a
distribution over plausible current values and decide whether to act or
defer. These systems should not necessarily be judged by the same scalar
recall metric.

This paper studies a stricter evaluation question: \textbf{when
long-horizon personalization drives action, which conclusions survive
changes in calibration protocol and evaluator assumptions?} We formulate
long-horizon personalization as partially observed state estimation
followed by risk-sensitive decision-making. PersonaMetrica-Bench
separates the observation layer from the belief-update layer so the
latter can be stress-tested with known provenance and scope. We then
evaluate both correctness and the relationship between confidence and
downstream action.

Our strongest result is not that one memory algorithm always wins.
Rather, the ranking depends on \textbf{both the metric and the
confidence protocol}. With native confidence, the 800-turn experiment
makes PBS look like a substantially better action substrate despite
slightly lower state accuracy. Once every tracker is given the same
development-only post-hoc calibration opportunity, the fixed-threshold
utility ordering flips and time-aware-latest becomes better at that
operating point. AURC, which depends on confidence ordering rather than
a particular action threshold, still favors PBS. This exposes two
evaluation confounds at once: state accuracy can miss selective-risk
structure, while fixed-threshold utility can reward one system merely
because its confidence scale is better matched to the chosen threshold.

\subsubsection{Contributions}\label{contributions}

\begin{enumerate}
\def\labelenumi{\arabic{enumi}.}
\tightlist
\item
  \textbf{Evaluation framing.} We formalize long-horizon personalization
  as belief-state estimation plus action under uncertainty and argue for
  reporting calibration and risk--coverage alongside state accuracy.
\item
  \textbf{PersonaMetrica-Bench.} A deterministic controlled testbed for
  protocol sensitivity with ten preference domains, durable changes,
  temporary exceptions, explicit/implicit evidence, noisy behavioral
  cues, hypothetical and other-person distractors, long-horizon
  variants, counterfactual interventions, and trajectory-level grader
  stress tests.
\item
  \textbf{Strong executable baselines.} Six memory/state systems ranging
  from first-mention and latest-observation to a time-aware baseline
  that filters distractors and handles temporary expiry.
\item
  \textbf{Personal Belief State.} A transparent evidence-weighted
  temporal representation with provenance, source reliability, recency,
  context scope, expiry, and probabilistic readout.
\item
  \textbf{Evaluation analysis.} Held-out multi-seed statistics,
  risk--coverage/AURC, a threshold/deferral-cost sensitivity sweep, and
  a development-only calibration protocol that reveals when
  native-confidence utility comparisons are not fair.
\item
  \textbf{Evaluator-as-system discipline.} Held-out grader attacks,
  causal trajectory contracts, abstention under missing evidence,
  implementation fingerprints, and disagreement routing test whether the
  measurement system itself generalizes. A 620-item human-annotation
  pack and vendor-neutral external-model evaluator are included as
  unexecuted infrastructure; no human or frontier-model result is
  fabricated.
\end{enumerate}

\subsubsection{Research questions}\label{research-questions}

\begin{itemize}
\tightlist
\item
  \textbf{RQ1 --- Metric sufficiency:} Does state accuracy rank
  personal-state systems the same way as selective risk and downstream
  action quality?
\item
  \textbf{RQ2 --- Calibration fairness:} Do conclusions from
  fixed-threshold utility survive when every system receives the same
  development-only confidence calibration opportunity?
\item
  \textbf{RQ3 --- Mechanism:} Which evidence-handling components of PBS
  account for its behavior under noise, temporary scope, and long
  histories?
\end{itemize}

The paper treats RQ2 as a first-class evaluation question rather than a
post-hoc implementation detail. Any utility comparison that gates action
on confidence is meaningful only relative to an explicit confidence
calibration protocol.

\subsection{2. Problem formulation}\label{problem-formulation}

Let a user have latent preference state \(z_t\). The agent observes
evidence \(x_t\), which may be explicit, implicit, noisy, hypothetical,
about another person, or scoped to a temporary context. A
personalization system maintains a belief

\[
b_t(z)=p(z_t\mid x_{1:t})
\]

and uses it to answer a query or choose an action \(a_t\).

A conventional state metric measures

\[
\mathrm{Acc}=\mathbb{E}[\mathbf{1}(\hat z_t=z_t)].
\]

For an acting assistant this is incomplete. We additionally measure:

\begin{itemize}
\tightlist
\item
  \textbf{coverage:} fraction of queries on which confidence is above an
  action threshold;
\item
  \textbf{selective accuracy:} accuracy conditional on action;
\item
  \textbf{Brier score:} calibration of confidence against correctness;
\item
  \textbf{risk--coverage curve:} error rate as progressively
  less-certain decisions are included;
\item
  \textbf{AURC:} area under the risk--coverage curve, lower is better;
\item
  \textbf{decision utility:} a transparent operating-point metric in
  which correct action is +1, wrong action is −1, and deferral has an
  explicit cost.
\end{itemize}

The utility coefficient is not assumed to be universal. We therefore
report a sensitivity sweep and treat AURC as the more defensible
threshold-free measure of confidence quality.

\subsection{3. Personal Belief State}\label{personal-belief-state}

PBS stores evidence rather than overwriting a single current value. Each
evidence record includes:

\begin{itemize}
\tightlist
\item
  semantic slot and candidate value;
\item
  turn and provenance;
\item
  source type: explicit, implicit choice, implicit behavior, inferred,
  hypothetical, or other-person;
\item
  observation confidence;
\item
  optional context scope;
\item
  optional temporary validity horizon.
\end{itemize}

At query time, each record receives a weight determined by source
reliability, observation confidence, exponential recency, context
compatibility, and temporary validity. Candidate values are normalized
to a probability distribution. Hypothetical and other-person evidence
remain in provenance but receive near-zero persistence weight. Temporary
evidence is strong in scope and sharply discounted after expiry rather
than erased.

PBS is intentionally transparent and dependency-light. It is not
presented as a replacement for a learned LLM memory module. Its role is
twofold: to instantiate the belief-state evaluation hypothesis and to
provide an auditable reference system whose failure modes can be studied
before introducing model-extraction errors.

\subsubsection{3.1 Development-only parameter
selection}\label{development-only-parameter-selection}

PBS has two continuous parameters in the checked-in experiment: evidence
half-life and softmax temperature. We select them by \textbf{minimum
mean Brier score on development seeds 1--3 only}, using 120 users per
seed. The selected values are a 90-turn half-life and temperature 0.20.
All reported ten-seed standard statistics use held-out seeds 10--19, and
long-horizon statistics use held-out seeds 30--39. The full search table
is released in \texttt{reports/BELIEF\_TUNING.md}.

\subsection{4. PersonaMetrica-Bench}\label{personametrica-bench}

\subsubsection{4.1 User state}\label{user-state}

The benchmark uses ten binary preference domains: work time, response
detail, travel pace, restaurant noise, notification style, budget style,
exercise time, meeting style, reading depth, and planning style.

\subsubsection{4.2 Evidence processes}\label{evidence-processes}

Each simulated user begins with explicit preferences and later receives
4--7 durable changes. Histories also contain:

\begin{itemize}
\tightlist
\item
  implicit-choice updates;
\item
  repeated behavioral cues that are \textbf{noisy} rather than always
  correct;
\item
  hypothetical distractors;
\item
  statements about another person;
\item
  temporary exceptions with explicit expiry;
\item
  queries that may fall inside or outside temporary scope.
\end{itemize}

The standard exported dataset contains \textbf{500 users, 27,133
structured observations, and 13,500 decision queries}. Surface-language
metadata spans English, Spanish, French, Persian, and Chinese. The
current benchmark begins \textbf{after} evidence extraction, so these
tags must not be interpreted as multilingual natural-language
understanding results.

\subsubsection{4.3 Distribution shifts}\label{distribution-shifts}

We evaluate five controlled regimes:

\begin{enumerate}
\def\labelenumi{\arabic{enumi}.}
\tightlist
\item
  \textbf{Standard:} 500 users × 240 turns.
\item
  \textbf{High noise:} more hypothetical/other-person distractors.
\item
  \textbf{Implicit-heavy:} more preference changes conveyed via implicit
  evidence.
\item
  \textbf{Temporary-heavy:} substantially more scoped exceptions.
\item
  \textbf{Long horizon:} 180 users × 800 turns in the checked-in shift
  run.
\end{enumerate}

\subsection{5. Baselines}\label{baselines}

We compare six executable systems:

\begin{itemize}
\tightlist
\item
  \textbf{First-mention:} stores the first observed value forever.
\item
  \textbf{Latest-observation:} blindly overwrites with the newest
  observation.
\item
  \textbf{Sliding-window-40:} uses the newest observation in a 40-turn
  window.
\item
  \textbf{Trusted-latest:} removes obvious hypothetical/other-person
  evidence but does not model temporary expiry or aggregate uncertainty.
\item
  \textbf{Time-aware-latest:} a strong symbolic baseline that filters
  distractors, distinguishes general and temporary evidence, and checks
  expiry at read time, but does not aggregate competing evidence.
\item
  \textbf{PBS:} evidence-weighted probabilistic belief state.
\end{itemize}

The inclusion of time-aware-latest is important: without it, the
benchmark would reward a method merely for knowing explicit scope
semantics. The repository records the earlier failure in which this
baseline initially solved a cleaner version of the benchmark, motivating
noisy behavioral evidence in the final generator.

These controlled baselines are not substitutes for model-level
comparisons with LongMemEval, PrefEval, HorizonBench, PerMemBench,
ProactiveBench, or production memory stacks. The release includes an
external-model contract specifically so those results can be added
without changing the metric implementation.

\subsection{6. Metrics and statistical
protocol}\label{metrics-and-statistical-protocol}

The primary controlled table reports state accuracy, action coverage,
selective accuracy, decision utility at threshold 0.70 and deferral cost
0.10, Brier score, and general/temporary slices.

For threshold-free confidence quality, we report the risk--coverage
curve and AURC. For sensitivity, thresholds \{0.50, 0.60, 0.70, 0.80,
0.90\} are crossed with deferral costs \{0, 0.05, 0.10, 0.20, 0.40,
0.60\}.

For the standard PBS comparison, ten held-out seeds (10--19) use 220
users × 240 turns each. For the long-horizon ranking analysis, ten
held-out seeds (30--39) use 90 users × 800 turns each. Bootstrap
intervals summarize \textbf{seed-level simulator variation}, not
uncertainty over a human population.

\subsection{7. Results}\label{results}

\subsubsection{7.1 Standard controlled
benchmark}\label{standard-controlled-benchmark}

On the checked-in standard run:

\begin{longtable}[]{@{}
  >{\raggedright\arraybackslash}p{(\columnwidth - 10\tabcolsep) * \real{0.1304}}
  >{\raggedleft\arraybackslash}p{(\columnwidth - 10\tabcolsep) * \real{0.1739}}
  >{\raggedleft\arraybackslash}p{(\columnwidth - 10\tabcolsep) * \real{0.1739}}
  >{\raggedleft\arraybackslash}p{(\columnwidth - 10\tabcolsep) * \real{0.1739}}
  >{\raggedleft\arraybackslash}p{(\columnwidth - 10\tabcolsep) * \real{0.1739}}
  >{\raggedleft\arraybackslash}p{(\columnwidth - 10\tabcolsep) * \real{0.1739}}@{}}
\toprule\noalign{}
\begin{minipage}[b]{\linewidth}\raggedright
Method
\end{minipage} & \begin{minipage}[b]{\linewidth}\raggedleft
State acc.
\end{minipage} & \begin{minipage}[b]{\linewidth}\raggedleft
Coverage
\end{minipage} & \begin{minipage}[b]{\linewidth}\raggedleft
Selective acc.
\end{minipage} & \begin{minipage}[b]{\linewidth}\raggedleft
Decision utility
\end{minipage} & \begin{minipage}[b]{\linewidth}\raggedleft
Brier ↓
\end{minipage} \\
\midrule\noalign{}
\endhead
\bottomrule\noalign{}
\endlastfoot
Trusted-latest & 0.853 & 0.631 & 0.908 & 0.478 & 0.129 \\
Time-aware-latest & 0.906 & 0.611 & \textbf{1.000} & 0.572 & 0.080 \\
\textbf{PBS} & \textbf{0.932} & \textbf{0.824} & 0.977 & \textbf{0.768}
& \textbf{0.060} \\
\end{longtable}

The strong time-aware baseline is perfectly accurate among the subset on
which it acts at threshold 0.70, but it acts on only 61.1\% of queries.
PBS covers 82.4\% while retaining 97.7\% selective accuracy.

Threshold-free AURC is \textbf{0.0107 for PBS}, compared with
\textbf{0.0208} for time-aware-latest and \textbf{0.0538} for
trusted-latest.

\begin{figure}
\centering
\includegraphics{paper/figures/accuracy_vs_utility.png}
\caption{Raw state accuracy and decision utility on the standard
controlled benchmark.}
\end{figure}

\subsubsection{7.2 Held-out standard
seeds}\label{held-out-standard-seeds}

PBS hyperparameters are frozen from development seeds 1--3. Across test
seeds 10--19, relative to time-aware-latest:

\begin{longtable}[]{@{}lrr@{}}
\toprule\noalign{}
PBS − time-aware-latest & Mean & 95\% bootstrap interval \\
\midrule\noalign{}
\endhead
\bottomrule\noalign{}
\endlastfoot
Decision utility & \textbf{+0.210} & \textbf{{[}+0.205, +0.216{]}} \\
State accuracy & \textbf{+0.024} & \textbf{{[}+0.021, +0.026{]}} \\
Selective accuracy & −0.020 & {[}−0.022, −0.018{]} \\
Brier reduction & \textbf{+0.020} & \textbf{{[}+0.019, +0.021{]}} \\
\end{longtable}

The selective-accuracy decrease is expected: PBS deliberately covers
many more cases rather than only high-confidence explicit observations.

\subsubsection{7.3 Long-horizon native-confidence
diagnostic}\label{long-horizon-native-confidence-diagnostic}

With each tracker using its \textbf{native, unmatched confidence scale},
800-turn histories produce an apparent ranking disagreement. In the
checked-in shift run, time-aware-latest has slightly \textbf{higher} raw
state accuracy than PBS (0.844 vs.~0.831) but lower fixed-threshold
decision utility (0.207 vs.~0.503). Across ten held-out long-horizon
seeds:

\begin{longtable}[]{@{}
  >{\raggedright\arraybackslash}p{(\columnwidth - 4\tabcolsep) * \real{0.2727}}
  >{\raggedleft\arraybackslash}p{(\columnwidth - 4\tabcolsep) * \real{0.3636}}
  >{\raggedleft\arraybackslash}p{(\columnwidth - 4\tabcolsep) * \real{0.3636}}@{}}
\toprule\noalign{}
\begin{minipage}[b]{\linewidth}\raggedright
PBS − time-aware-latest
\end{minipage} & \begin{minipage}[b]{\linewidth}\raggedleft
Mean
\end{minipage} & \begin{minipage}[b]{\linewidth}\raggedleft
95\% bootstrap interval
\end{minipage} \\
\midrule\noalign{}
\endhead
\bottomrule\noalign{}
\endlastfoot
State accuracy & \textbf{−0.014} & \textbf{{[}−0.017, −0.012{]}} \\
Native-confidence decision utility & \textbf{+0.284} &
\textbf{{[}+0.276, +0.291{]}} \\
Coverage & \textbf{+0.328} & \textbf{{[}+0.322, +0.333{]}} \\
Brier reduction & \textbf{+0.016} & \textbf{{[}+0.014, +0.017{]}} \\
\end{longtable}

This comparison is intentionally treated as a \textbf{diagnostic, not a
fair model-selection result}, because the systems' confidence scales
have not yet been equalized. Section 7.6 applies the same
development-only calibration opportunity to every tracker and shows that
the fixed-threshold utility ordering reverses. The stable conclusion is
therefore about protocol sensitivity, not a universal PBS utility
advantage.

\begin{figure}
\centering
\includegraphics{paper/figures/risk_coverage.png}
\caption{Risk--coverage curve for the strongest controlled systems.}
\end{figure}

\subsubsection{7.4 Distribution shifts}\label{distribution-shifts-1}

PBS has the highest checked-in decision utility in all five regimes. Its
utility remains 0.736 under high noise, 0.705 under implicit-heavy
histories, 0.773 under temporary-heavy histories, and 0.503 at 800
turns.

\begin{figure}
\centering
\includegraphics{paper/figures/distribution_shifts.png}
\caption{Decision utility under controlled distribution shifts.}
\end{figure}

\subsubsection{7.5 Additional robustness
analyses}\label{additional-robustness-analyses}

Sensitivity sweeps, component ablations, a 48-cell stress grid, and
exact paired tests are reported in the technical appendix and released
artifacts. Briefly, PBS is preferred to time-aware-latest in
\textbf{23/30} threshold/cost cells before shared calibration;
provenance reliability and evidence aggregation are the largest PBS
ablation effects; and PBS has higher utility in \textbf{46/48}
controlled stress cells but is not uniformly superior. These analyses
narrow rather than broaden the claim: operating-point conclusions depend
on calibration and cost assumptions, while mechanism effects are
generator-dependent. A supplementary trajectory-level agent-evaluation
extension also tests whether graders distinguish successful final states
from policy-compliant, verified trajectories; it is reported in Appendix
G and is not used as a frontier-model capability claim. Appendix I
further evaluates the frozen grader on five held-out attack families,
and Appendix J adds paired counterfactual checks of belief-state
robustness versus adaptivity.

\subsubsection{7.6 Equal-opportunity confidence calibration changes the
conclusion}\label{equal-opportunity-confidence-calibration-changes-the-conclusion}

A fixed action threshold is only comparable when confidence scales are
comparable. We therefore fit the same deliberately low-capacity post-hoc
calibrator for \textbf{every tracker} on development seeds 1--3 only:
each raw-confidence bucket is mapped to Laplace-smoothed empirical
correctness, then frozen for held-out evaluation. No test labels are
used.

On 800-turn held-out seeds, this changes the conclusion. The calibrated
time-aware-latest baseline obtains \textbf{0.687 decision utility},
compared with \textbf{0.588 for PBS}. PBS still has higher selective
accuracy (\textbf{0.908 vs.~0.844}) and slightly lower Brier
(\textbf{0.119 vs.~0.122}), but its lower coverage (\textbf{0.750
vs.~1.000}) is penalized at the chosen deferral cost. Thus the earlier
native-confidence utility advantage (\textbf{+0.284} on average) is not
protocol-invariant; after shared calibration the sign reverses to
approximately \textbf{−0.099}.

This is a central negative result. It does \textbf{not} imply that
time-aware-latest is globally better, nor that PBS is globally better.
It implies that fixed-threshold utility combines three
choices---confidence scale, action threshold, and deferral cost---and
can therefore change system ordering when those choices change. AURC is
less exposed to this particular confound because monotone post-hoc
calibration does not alter confidence ordering; in the checked-in
benchmark AURC remains \textbf{0.0107 for PBS vs.~0.0208 for
time-aware-latest}.

We therefore recommend a reporting hierarchy: (1) state accuracy, (2)
calibration error, (3) risk--coverage/AURC, and (4) operating-point
utility only after an explicitly shared calibration protocol. The
complete mappings and held-out results are released in
\texttt{reports/CALIBRATED\_BASELINES.md} and
\texttt{reports/CALIBRATION\_PROTOCOL\_ANALYSIS.md}.

A broader frozen protocol sweep reaches the same conclusion without
relying on a single threshold or calibrator. The focused
PBS-vs-time-aware slice reverses frequently across 75 cells, and the
full audit ranks \textbf{all six released baselines} under the same
protocol family. Time-aware-latest wins 45 cells, PBS 27, and
first-mention 3. Across unordered pairs of protocol cells, the
probability of selecting a different winner is \textbf{0.516}; mean
complete-ranking Kendall \(\tau\) is \textbf{0.699}, with a minimum of
\textbf{0.20}. To test whether this headline is a seed or grid artifact,
a separate 10-user-per-seed diagnostic runs \textbf{2,000 paired seed
bootstraps}, yielding winner-change probability \textbf{0.516 {[}0.507,
0.538{]}} and mean Kendall \(\tau\) \textbf{0.702 {[}0.692, 0.720{]}}.
Thirteen leave-one-level-out grid perturbations all retain multiple
winners; winner-change probability ranges \textbf{0.327--0.564}.
Appendix K reports the focused matrix, Appendix L the full-leaderboard
and benchmark-health audit, and Appendix M the seed/grid robustness
checks. These intervals quantify synthetic-generator seed variation
only; they are not population estimates or external leaderboard results.

\subsection{8. Relation to prior work}\label{relation-to-prior-work}

The personalization benchmark landscape is now crowded, so
PersonaMetrica does \textbf{not} claim to introduce the first benchmark
for preference following, evolving preferences, personalized memory, or
proactive assistance. PrefEval evaluates explicit and implicit
preference following in long-context conversations (S. Zhao et al.
2025). PersonaLens evaluates personalized task-oriented assistants with
simulated users and LLM judges (Z. Zhao et al. 2025). HorizonBench
directly studies evolving preferences over six-month synthetic histories
and diagnoses stale-belief failures (Li et al. 2026). BeliefShift
evaluates temporal belief consistency, contradiction resolution, and
evidence-driven revision (Myakala, Agrawal, and Manche 2026). PERMA
models event-driven preference evolution in realistic task environments
(Liu et al. 2026). PerMemBench studies personalized memory-storage
policies (In et al. 2026), while PAHF studies continual personalization
with pre-action clarification and post-action feedback (Liang et al.
2026). pi-Bench evaluates proactive personal assistants with hidden
intents, tools, and persistent sessions (Zhang et al. 2026). OP-Bench
targets over-personalization such as irrelevance, repetition, and
sycophancy (Hu et al. 2026).

PersonaMetrica targets a different layer: \textbf{measurement validity
when personalization systems emit confidence and that confidence gates
action}. Its core empirical object is not a new memory leaderboard but
\emph{ranking stability under defensible evaluation changes}.
Native-confidence utility and shared-calibration utility can select
different systems at the same nominal action threshold, while
threshold-free risk--coverage tells a different story again. The
trajectory extension then applies the same skepticism to the evaluator:
a grader that is perfect on authored perturbations can fail
catastrophically on held-out attack families, motivating causal-contract
checks, \texttt{Unknown} under missing evidence, grader fingerprints,
and disagreement review.

This positioning is deliberately narrower than a ``first'' claim.
Related work in selective prediction already establishes risk--coverage
and abstention as general evaluation tools, while 2026 work studies
ranking instability under prompt/evaluator variation (Du et al. 2026;
Mahmood, Abdaljalil, and Kurban 2026), disagreement-aware stable
aggregation (Bonagiri et al. 2026), long-horizon trajectory attribution
(Chen et al. 2026), and trajectory-level safety calibration under
adaptive attacks (Asif et al. 2026). PersonaMetrica therefore does
\textbf{not} claim that rank instability, trajectory attacks, or grader
auditing are individually new. Its narrower contribution is to connect
protocol-sensitive system selection with \textbf{evolving personal
state}, confidence-gated action, counterfactual preference
interventions, and evaluator validity in one reproducible artifact.

\subsection{9. Human and model evaluation
protocol}\label{human-and-model-evaluation-protocol}

The release contains \textbf{620 independently labelable tasks} (300
proactivity decisions, 320 belief-update decisions) with a browser
interface that hides simulator reference labels, plus a provider-neutral
external-model evaluator reporting accuracy, Brier, and ECE. It also
includes an executable adapter for the official HorizonBench per-item
result JSONL contract; the checked-in adapter run uses synthetic schema
fixtures only. These are \textbf{unexecuted external-evidence
infrastructure}: no human or frontier-model result is reported. The
intended next study holds the evaluation protocol fixed while replacing
simulator-derived evidence or labels with independent human/model
outputs.

\subsection{10. Limitations}\label{limitations}

PersonaMetrica-Bench begins \textbf{after structured evidence
extraction}, uses synthetic users, binary preference domains, designed
source reliabilities, and generator-defined transition processes. PBS
parameters are selected on development simulator seeds and should not be
interpreted as learned human preference dynamics. The benchmark can
exhibit generator--method coupling, and multilingual tags are metadata
rather than multilingual understanding results. No frontier-model table
or human-subject study is reported. Decision utility is task-dependent,
and this paper directly shows that its ranking can reverse after shared
confidence calibration; utility must therefore be interpreted only with
calibration, threshold, and deferral cost specified. Finally, the
current Croissant metadata and anonymous bundle are complete locally,
but venue-compliant external anonymous hosting remains a submission
requirement.

\subsection{11. Reproducibility and responsible
release}\label{reproducibility-and-responsible-release}

All controlled primary experiments run on CPU; the release records
seeds, generated data, scripts, figures, compute measurements, integrity
audits, a dataset card, Croissant metadata, and an anonymous submission
bundle. The benchmark contains only synthetic users. Because
personal-state systems can enable persistent profiling or overconfident
action, the broader system artifact includes forgetting, privacy-policy
hooks, permission boundaries, provenance, and outcome verification. No
human-subject result supports the present claims.

The practical evaluation implication is simple: when personal state can
trigger action, report \textbf{state accuracy, calibration,
risk--coverage/AURC, and any operating-point utility together}, and
apply a shared confidence calibration protocol before comparing
fixed-threshold action policies.

\subsection{12. Conclusion}\label{conclusion}

Personalization is a belief-update problem before it is a retrieval
problem, and an action problem after it. PersonaMetrica-Bench shows that
system rankings can change not only when moving from state accuracy to
action metrics, but also when the \textbf{confidence calibration and
action protocol itself} changes. In the controlled long-horizon setting,
native-confidence utility favors PBS, whereas equal development-only
post-hoc calibration can favor time-aware-latest at the same operating
point; threshold-free AURC still favors PBS. Across the full six-system
75-cell audit, even the leaderboard winner and complete ordering vary
substantially. The larger contribution is therefore evaluative rather
than algorithmic: long-horizon personal assistants should be compared
with explicit calibration protocols, threshold-free risk metrics, and
ranking-stability diagnostics rather than a single operating-point
leaderboard. PBS is one transparent reference implementation for
studying those interactions, not a claim of a universally superior
memory architecture.

\section*{References}\label{bibliography}
\addcontentsline{toc}{section}{References}

\phantomsection\label{refs}
\begin{CSLReferences}{1}{0}
\bibitem[\citeproctext]{ref-asif2026blindspot}
Asif, Sadia, Mohammad Mohammadi Amiri, Momin Abbas, Tejaswini Pedapati,
and Prasanna Sattigeri. 2026. {``BLINDSPOT: A Benchmark for Safety and
Refusal Calibration in Long-Horizon Tool-Using Agents.''} \emph{arXiv
Preprint arXiv:2609.16305}.

\bibitem[\citeproctext]{ref-bonagiri2026stableval}
Bonagiri, Akash, Gerard Janno Anderias, Saee Patil, Angelina Lai, Devang
Borkar, Gezheng Kang, Ishant Gandhi, Setareh Rafatirad, and Houman
Homayoun. 2026. {``STABLEVAL: Disagreement-Aware and Stable Evaluation
of AI Systems.''} \emph{arXiv Preprint arXiv:2605.02122}.

\bibitem[\citeproctext]{ref-chen2026trajectoryattribution}
Chen, Jing, Yang Sun, Li Zhang, Lin Xu, and Jie Shi. 2026.
{``Long-Horizon Agent Trajectory Attribution: A Unified Benchmark and
Fine-Grained Annotation Framework.''} \emph{arXiv Preprint
arXiv:2608.06909}.

\bibitem[\citeproctext]{ref-du2026promptstability}
Du, Shaoshuai, Penghao Liang, Yixian Shen, Chuanqi Shi, Hang Zhang, and
Lun Wang. 2026. {``On the Stability of Prompt Ranking in Large Language
Model Evaluation.''} \emph{arXiv Preprint arXiv:2606.24381}.

\bibitem[\citeproctext]{ref-hu2026opbench}
Hu, Yulin, Zimo Long, Jiahe Guo, Xingyu Sui, Xing Fu, Weixiang Zhao,
Yanyan Zhao, and Bing Qin. 2026. {``OP-Bench: Benchmarking
over-Personalization for Memory-Augmented Personalized Conversational
Agents.''} \emph{arXiv Preprint arXiv:2601.13722}.

\bibitem[\citeproctext]{ref-in2026permem}
In, Yeonjun, Wonjoong Kim, Sangwu Park, Kanghoon Yoon, and Chanyoung
Park. 2026. {``Personalize-Then-Store: Benchmarking and Learning
Personalized Memory for Long-Horizon Agents.''} \emph{arXiv Preprint
arXiv:2605.25535}.

\bibitem[\citeproctext]{ref-li2026horizonbench}
Li, Shuyue Stella, Bhargavi Paranjape, Kerem Oktar, Zhongyao Ma, Gelin
Zhou, Lin Guan, Na Zhang, et al. 2026. {``HorizonBench: Long-Horizon
Personalization with Evolving Preferences.''} \emph{arXiv Preprint
arXiv:2604.17283}.

\bibitem[\citeproctext]{ref-liang2026pahf}
Liang, Kaiqu, Julia Kruk, Shengyi Qian, Xianjun Yang, Shengjie Bi,
Yuanshun Yao, Shaoliang Nie, et al. 2026. {``Learning Personalized
Agents from Human Feedback.''} \emph{arXiv Preprint arXiv:2602.16173}.

\bibitem[\citeproctext]{ref-perma2026}
Liu, Shuochen, Junyi Zhu, Long Shu, Junda Lin, Yuhao Chen, Haotian
Zhang, Chao Zhang, et al. 2026. {``PERMA: Benchmarking Personalized
Memory Agents via Event-Driven Preference and Realistic Task
Environments.''} \emph{arXiv Preprint arXiv:2603.23231}.

\bibitem[\citeproctext]{ref-mahmood2026rankreversal}
Mahmood, Alhasan, Samir Abdaljalil, and Hasan Kurban. 2026. {``Rank
Reversal in Multilingual LLM Judges: A Label-Free Double-Centering
Calibrator.''} \emph{arXiv Preprint arXiv:2608.22432}.

\bibitem[\citeproctext]{ref-myakala2026beliefshift}
Myakala, Praveen Kumar, Manan Agrawal, and Rahul Manche. 2026.
{``BeliefShift: Benchmarking Temporal Belief Consistency and Opinion
Drift in LLM Agents.''} \emph{arXiv Preprint arXiv:2603.23848}.

\bibitem[\citeproctext]{ref-zhang2026pibench}
Zhang, Haoran, Luxin Xu, Zhilin Wang, Runquan Gui, Shunkai Zhang, Haodi
Lei, Zihao He, et al. 2026. {``\(\\pi\)-Bench: Evaluating Proactive
Personal Assistant Agents in Long-Horizon Workflows.''} \emph{arXiv
Preprint arXiv:2605.14678}.

\bibitem[\citeproctext]{ref-zhao2025prefeval}
Zhao, Siyan, Mingyi Hong, Yang Liu, Devamanyu Hazarika, and Kaixiang
Lin. 2025. {``Do LLMs Recognize Your Preferences? Evaluating
Personalized Preference Following in LLMs.''} In \emph{International
Conference on Learning Representations}.

\bibitem[\citeproctext]{ref-zhao2025personalens}
Zhao, Zheng, Clara Vania, Subhradeep Kayal, Naila Khan, Shay B. Cohen,
and Emine Yilmaz. 2025. {``PersonaLens: A Benchmark for Personalization
Evaluation in Conversational AI Assistants.''} In \emph{Findings of the
Association for Computational Linguistics: ACL 2025}.
\url{https://doi.org/10.18653/v1/2025.findings-acl.927}.

\end{CSLReferences}

\section{Technical appendix}\label{technical-appendix}

\subsection{Appendix A. Benchmark integrity
audit}\label{appendix-a.-benchmark-integrity-audit}

The standard generator was independently re-audited from source in
v4.1.0. It deterministically produces \textbf{27,133 observations and
13,500 queries} for 500 users at seed 31. Replaying the first 25 users
produces an identical SHA-256 digest. Exact structured fingerprints
repeat across users because the benchmark intentionally uses a finite
schema; the scientific unit is the evolving user history, not textual
novelty. Full source/query distributions are in
\texttt{reports/BENCHMARK\_INTEGRITY\_AUDIT.md}.

\subsection{Appendix B. Experimental
details}\label{appendix-b.-experimental-details}

Development seeds are 1--3; standard held-out test seeds are 10--19;
long-horizon held-out seeds are 30--39. PBS half-life and temperature
are selected only by development Brier score. The action threshold used
in the main operating-point table is 0.70 and deferral cost is 0.10; the
sensitivity grid reports thresholds 0.50--0.90 and deferral costs
0--0.60. All component ablations reuse frozen PBS hyperparameters.

\subsection{Appendix C. Exact metric
definitions}\label{appendix-c.-exact-metric-definitions}

For query \(i\), let \(y_i\in\{0,1\}\) indicate whether the predicted
current state is correct and let \(c_i\in[0,1]\) be the system
confidence assigned to that prediction. For an action threshold
\(\tau\), the system acts when \(c_i\ge\tau\) and otherwise defers. We
report:

\begin{itemize}
\tightlist
\item
  \textbf{State accuracy:} \(N^{-1}\sum_i y_i\).
\item
  \textbf{Coverage:} \(N^{-1}\sum_i \mathbf{1}[c_i\ge\tau]\).
\item
  \textbf{Selective accuracy:}
  \(\sum_i y_i\mathbf{1}[c_i\ge\tau] / \sum_i \mathbf{1}[c_i\ge\tau]\),
  defined as zero when coverage is zero.
\item
  \textbf{Brier score:} \(N^{-1}\sum_i(c_i-y_i)^2\).
\item
  \textbf{Decision utility:} correct action \(+1\), incorrect action
  \(-1\), and deferral \(-d\), averaged over all queries, where the main
  operating point uses \(d=0.10\).
\item
  \textbf{Risk--coverage curve:} for each target coverage, queries are
  sorted by decreasing confidence and risk is \(1-\)accuracy on the
  retained prefix. \textbf{AURC} is the numerical area under this curve;
  lower is better.
\end{itemize}

The paper treats state accuracy, Brier, and AURC as protocol-level
metrics. Fixed-threshold utility is secondary because it additionally
depends on confidence calibration, \(\tau\), and \(d\).

\subsection{Appendix D. Equal-opportunity confidence
calibration}\label{appendix-d.-equal-opportunity-confidence-calibration}

To test whether native confidence gives some trackers an unfair
advantage, every tracker receives the same development-only post-hoc
calibration opportunity. Raw confidence values are rounded to two
decimals; for each value, development correctness is converted to a
Laplace-smoothed empirical probability \((k+1)/(n+2)\). The mapping is
fit on development seeds 1--3 only and frozen before standard held-out
seeds 10--19 and long-horizon seeds 30--39. Unseen raw-confidence values
map to the nearest development bucket. This calibrator is intentionally
low capacity: it removes the easiest confidence-scale mismatch without
introducing a learned model.

The long-horizon utility ordering changes after this calibration: native
confidence gives PBS a mean +0.284 utility advantage over
time-aware-latest, while shared development calibration yields
approximately -0.099. The paper therefore does not use the
native-confidence utility ranking as evidence of universal method
superiority.

\subsection{Appendix E. Claim-to-evidence
matrix}\label{appendix-e.-claim-to-evidence-matrix}

\begin{longtable}[]{@{}
  >{\raggedright\arraybackslash}p{(\columnwidth - 4\tabcolsep) * \real{0.3333}}
  >{\raggedright\arraybackslash}p{(\columnwidth - 4\tabcolsep) * \real{0.3333}}
  >{\raggedright\arraybackslash}p{(\columnwidth - 4\tabcolsep) * \real{0.3333}}@{}}
\toprule\noalign{}
\begin{minipage}[b]{\linewidth}\raggedright
Claim
\end{minipage} & \begin{minipage}[b]{\linewidth}\raggedright
Primary evidence
\end{minipage} & \begin{minipage}[b]{\linewidth}\raggedright
Important exclusion
\end{minipage} \\
\midrule\noalign{}
\endhead
\bottomrule\noalign{}
\endlastfoot
State accuracy alone is incomplete for acting assistants & standard +
800-turn comparisons; risk--coverage & synthetic structured evidence
only \\
Fixed-threshold utility is calibration-sensitive &
development-calibrated baseline analysis & simple bucket calibrator; no
frontier models \\
PBS confidence ordering is stronger in the checked-in benchmark & AURC
0.0107 vs.~0.0208 & does not imply better utility at every operating
point \\
PBS behavior depends strongly on provenance/aggregation & held-out
component ablations & effect sizes are generator-dependent \\
Findings persist across many controlled regimes & 48-cell stress grid &
two cells favor the baseline \\
\end{longtable}

\subsection{Appendix F. Additional robustness analyses moved from the
main
text}\label{appendix-f.-additional-robustness-analyses-moved-from-the-main-text}

\subsubsection{7.5 Sensitivity to action threshold and deferral
cost}\label{sensitivity-to-action-threshold-and-deferral-cost}

At the tested 180-user sensitivity seed, PBS has higher utility than
trusted-latest in \textbf{30/30} threshold/cost cells and higher utility
than time-aware-latest in \textbf{23/30} cells. The remaining cells
matter: utility is not a property of the memory system alone but of the
decision rule and the cost assigned to deferral. For this reason, the
paper treats risk--coverage/AURC as primary and single-threshold utility
as an operating-point analysis.

\subsubsection{7.6 Which PBS components
matter?}\label{which-pbs-components-matter}

We ablate PBS components on held-out seeds 10--19 (160 users × 240 turns
per seed), without retuning the remaining components. Removing recency
lowers decision utility by \textbf{20.7 pp}; replacing provenance-aware
source reliability with uniform reliability lowers it by \textbf{71.0
pp}; and removing evidence aggregation lowers it by \textbf{74.6 pp}.
Temporary expiry and context scope have smaller but measurable effects
on this generator.

\begin{longtable}[]{@{}lrrrr@{}}
\toprule\noalign{}
Variant & State acc. & Coverage & Decision utility & Brier ↓ \\
\midrule\noalign{}
\endhead
\bottomrule\noalign{}
\endlastfoot
Full PBS & 0.937 & 0.831 & \textbf{0.779} & \textbf{0.058} \\
No recency & 0.787 & 0.858 & 0.572 & 0.149 \\
No provenance reliability & 0.537 & 0.830 & 0.070 & 0.376 \\
No temporary expiry & 0.932 & 0.821 & 0.765 & 0.061 \\
No context scope & 0.928 & 0.829 & 0.761 & 0.065 \\
No evidence aggregation & 0.479 & 0.121 & 0.033 & 0.150 \\
\end{longtable}

\begin{figure}
\centering
\includegraphics{paper/figures/pbs_component_ablations.png}
\caption{PBS component ablations on held-out simulator seeds.}
\end{figure}

The large provenance and aggregation effects are also a warning about
benchmark-method coupling: PersonaMetrica-Bench explicitly varies
evidence source and contradiction structure, so methods that ignore
those fields are expected to degrade. The benchmark therefore treats
these ablations as mechanism diagnostics, not evidence that the same
effect sizes will hold on natural conversations.

\subsubsection{7.7 Stress grid and failure
region}\label{stress-grid-and-failure-region}

We additionally sweep \textbf{48 controlled cells}: four
distractor-noise rates × four behavioral-evidence error rates × three
temporary-preference rates, with 90 users × 240 turns per cell. PBS has
higher decision utility than time-aware-latest in \textbf{46/48} cells
and higher raw state accuracy in \textbf{34/48}. Mean utility gain is
\textbf{+0.138}, but the worst cell is \textbf{−0.025}. This failure
region matters: PBS is not uniformly superior when source reliabilities
are strongly misspecified.

\begin{figure}
\centering
\includegraphics{paper/figures/belief_stress_grid.png}
\caption{Stress grid averaged over distractor-noise levels.}
\end{figure}

This result narrows the claim: PBS is a useful reference implementation
for the proposed evaluation framing, not a universally optimal state
estimator.

\subsubsection{7.8 Paired held-out tests}\label{paired-held-out-tests}

The primary ten-seed comparisons are paired by simulator seed. Exact
two-sided sign tests and exact sign-flip permutation tests give
\textbf{p = 0.00195} for the standard utility gain, standard accuracy
gain, long-horizon utility gain, long-horizon accuracy reversal, and
long-horizon Brier reduction. All ten long-horizon seeds show higher PBS
utility while all ten show lower PBS raw state accuracy. These tests
quantify consistency across the pre-specified simulator seeds; they are
\textbf{not} inferential claims about a human population. More
importantly, these native-confidence utility tests do not by themselves
establish a fair action comparison across systems; Section 7.6 repeats
the comparison after a shared development-only calibration protocol.

The corresponding machine-readable outputs are
\texttt{data/belief\_sensitivity.json},
\texttt{data/pbs\_component\_ablations.json},
\texttt{data/belief\_stress\_grid.json}, and
\texttt{data/belief\_paired\_tests.json}.

\subsection{Appendix G. Trajectory-aware agent evaluation
extension}\label{appendix-g.-trajectory-aware-agent-evaluation-extension}

The release includes a controlled agent-evaluation extension using the
task/trial/grader/trajectory decomposition. \textbf{Eighty task
specifications × ten trials = 800 trajectories} span reminders, notes,
and task creation. Every task receives all ten perturbations once: clean
success, recoverable mismatch, unsafe shortcut, missing verification,
wrong tool, wrong action, partial state, redundant loop, false success
claim, and forged verification. The suite is diagnostic rather than a
frontier-model benchmark.

An end-state-only grader achieves \textbf{68.3\% accuracy} and falsely
accepts \textbf{44.3\%} of gold-failing trials. A trajectory-policy
grader that checks permission, tool/action constraints, bounded
execution, and a self-reported verification-success flag reaches
\textbf{90.0\% accuracy} but still has \textbf{13.9\% false accepts}.
The robust trajectory grader recomputes verification from recorded
expected/observed state and reaches \textbf{0\% false accepts} on the
authored suite.

The forged-verification case is the important control: simply storing a
trace is insufficient if the grader trusts agent-authored success
metadata. The evaluation needs independently checkable state evidence.
Full results are in \texttt{reports/AGENT\_TRAJECTORY\_EVAL.md} and
\texttt{data/agent\_trajectory\_eval.json}.

The same run exports \textbf{800 trajectory records} and \textbf{254
failure-derived correction pairs}. Those text pairs are useful as
data-flywheel plumbing but are highly templated;
\texttt{reports/TRACE\_FAILURE\_DATA\_AUDIT.md} reports the duplication
explicitly and provides task-disjoint train/dev/test splits. Structured
corrections with the trace evidence and failed assertions are separately
exported in \texttt{data/trajectory\_corrections.jsonl}. No
post-training generalization claim is made from these artifacts.

\subsection{Appendix H. Grader reliability, drift, and review
routing}\label{appendix-h.-grader-reliability-drift-and-review-routing}

The evaluator itself is stress-tested rather than assumed correct. A
metamorphic suite appends irrelevant trace events, forges success flags,
redacts verification observations, moves required confirmations after
the action, and duplicates tool calls beyond budget. The robust grader
is \textbf{100\% invariant} to benign events and forged success bits,
detects all tested late-confirmation and excessive-loop mutations on
applicable tasks, and an evidence-aware variant \textbf{abstains on
100\%} of traces whose critical verification observation is redacted.

A weighted composite grader also exhibits threshold drift: on the
checked-in controlled suite, lowering the pass threshold from 0.95 to
0.55 raises false accepts from roughly \textbf{14\% to 58\%}. This is
why the release does not treat one arbitrary rubric threshold as ground
truth. Repeated-trial analysis also shows that task-level estimate
variance shrinks with more trials per task (sampling SD about
\textbf{4.1 pp with one trial vs.~1.5 pp with eight trials} in the
checked-in reliability run).

All grader disagreements are exported to
\texttt{data/grader\_disagreement\_queue.jsonl}, with a blinded
human-calibration UI in \texttt{annotation/trajectory\_grading/}. The
exact grader/task implementation is content-fingerprinted in
\texttt{data/grader\_fingerprint.json}; evaluator-code changes require
regenerated results. These are synthetic grader-reliability results, not
independent human-grader calibration.

\subsection{Appendix I. Held-out evaluator
attacks}\label{appendix-i.-held-out-evaluator-attacks}

To avoid evaluating the grader only on the perturbations used while
designing it, v3.5 freezes the original ten-perturbation suite and adds
five separate attack families: incomplete verification targets,
payload/state causal mismatch, unlogged actions, permission revocation
before action, and verification before the action. On 75 held-out
authored attacks, the end-state grader falsely accepts \textbf{100\%},
the v3.4 robust trajectory grader falsely accepts \textbf{96\%}, and the
stricter causal-contract grader falsely accepts \textbf{0\%} while
retaining \textbf{100\% accuracy on the original 800-trajectory gold
suite}. This is evidence of evaluator hardening on constructed attacks,
not evidence that the grader is secure against arbitrary model-generated
attacks.

\subsection{Appendix J. Counterfactual belief-state
checks}\label{appendix-j.-counterfactual-belief-state-checks}

We additionally run 540 paired structured-evidence interventions.
Time-aware-latest and PBS are both highly sensitive to an explicit
current update (100\% and \textbf{99.3\%}, respectively).
Time-aware-latest flips on weak conflicting behavioral evidence in
\textbf{90.0\%} of pairs, whereas PBS flips in \textbf{35.6\%}. PBS is
invariant to a high-confidence other-person distractor in
\textbf{98.0\%} of pairs, leaving a small but concrete \textbf{2.0\%
nuisance-sensitivity failure rate}. These diagnostics expose a
robustness/adaptivity tradeoff and are not natural-language
understanding results.

\subsection{Appendix K. Protocol Stability
Matrix}\label{appendix-k.-protocol-stability-matrix}

To test whether the central rank-reversal result depends on a single
hand-picked operating point, we freeze a broader protocol grid before
interpreting the long-horizon diagnostic cohort: three confidence
treatments (native, development-only bucket calibration,
development-only isotonic calibration), five action thresholds, and five
deferral costs, for \textbf{75 protocol cells}. Calibrators are fit on
development seeds 1--3 only. The diagnostic cohort uses 40 users for
each of long-horizon seeds 30--39; the larger primary benchmark remains
unchanged.

PBS has higher mean utility in \textbf{30/75} cells and
time-aware-latest in \textbf{45/75}. The probability that two distinct
protocol cells select opposite pairwise winners is \textbf{0.486};
winner entropy is \textbf{0.971 bits}. Native confidence favors PBS in
20/25 cells, whereas each development-calibrated family favors
time-aware-latest in 20/25 cells. This is a measurement-fragility
result, not a claim that one calibration family is inherently correct or
that the same reversal frequency holds for frontier models.

\begin{figure}
\centering
\includegraphics{paper/figures/protocol_stability_matrix.png}
\caption{Protocol Stability Matrix. Values are mean PBS minus
time-aware-latest utility across long-horizon diagnostic seeds; sign
changes indicate protocol-dependent winner changes.}
\end{figure}

The exact grid is stored in \texttt{configs/protocol\_registry.json},
fingerprinted in \texttt{data/protocol\_registry\_fingerprint.json}, and
validated by \texttt{scripts/validate\_protocol\_registry.py}. This is a
\textbf{release-level protocol lock}, not an externally timestamped
preregistration. Full results are in
\texttt{reports/PROTOCOL\_STABILITY\_MATRIX.md} and
\texttt{data/protocol\_stability\_matrix.json}.

\subsection{Appendix L. Full-Leaderboard Stability and Benchmark
Health}\label{appendix-l.-full-leaderboard-stability-and-benchmark-health}

The focused pairwise protocol sweep can understate evaluation
instability because a third method may become competitive under a
different confidence/cost regime. We therefore replay \textbf{all six
released baselines} on the same long-horizon diagnostic histories across
the frozen 75-cell protocol family. Calibration is fit on development
seeds 1-3 and frozen before long-horizon seeds 30-39.

\begin{longtable}[]{@{}lr@{}}
\toprule\noalign{}
Full-leaderboard quantity & Value \\
\midrule\noalign{}
\endhead
\bottomrule\noalign{}
\endlastfoot
Time-aware-latest winner cells & 45 / 75 \\
PBS winner cells & 27 / 75 \\
First-mention winner cells & 3 / 75 \\
Probability two protocol cells choose different winners &
\textbf{0.516} \\
Mean Kendall \(\tau\) across complete rankings & \textbf{0.699} \\
Minimum Kendall \(\tau\) & \textbf{0.20} \\
\end{longtable}

The three first-mention wins are treated as a \textbf{warning}, not a
positive result for that baseline: an operating protocol can reward
coverage/confidence economics strongly enough that a stale estimator
becomes the apparent winner. This is why PersonaMetrica reports
threshold-free state/risk metrics alongside operating-point utility.

Benchmark health is also release-gated. Synthetic identities are
seed-scoped; development and test display IDs have zero overlap; after
identifiers are removed, development and test complete-history
fingerprints have zero overlap; the 500-user standard cohort has zero
duplicate complete histories; all ten preference domains are queried;
temporary-scope queries are present; answer skew remains bounded; and
the benchmark is not saturated. Structured event fragments can repeat
because the simulator uses a finite schema, but the evolving
full-history unit does not duplicate in the checked-in standard cohort.

\subsection{Appendix M. Seed uncertainty, grid robustness, and
external-result
contract}\label{appendix-m.-seed-uncertainty-grid-robustness-and-external-result-contract}

The full-leaderboard result is additionally stress-tested along two axes
that were not used to choose the reported winner counts. First, a
lower-cost diagnostic cohort with 10 users for each long-horizon seed
replays all six methods across the same 75 protocol cells. We resample
the ten long-horizon seeds \textbf{2,000 times as paired units}, using
identical resampled seed indices for every method and protocol.
Winner-change probability is \textbf{0.516} with a 95\% seed-bootstrap
interval of \textbf{{[}0.507, 0.538{]}}; mean complete-ranking Kendall
\(\tau\) is \textbf{0.702 {[}0.692, 0.720{]}}. In all bootstrap
replicates, winner-change probability remains above 0.40. The least
stable individual protocol-cell winner has bootstrap support
\textbf{0.619}, so the release does not treat every cell winner as
certain. This diagnostic uses fewer users per seed than the primary
40-user-per-seed leaderboard audit and is reported as robustness
evidence, not a replacement primary estimate.

Second, we run \textbf{13 leave-one-level-out perturbations} of the
frozen protocol grid: each calibration family, action threshold, and
deferral cost is omitted once without inventing replacement levels from
test outcomes. Every perturbed grid retains more than one winner.
Winner-change probability ranges \textbf{0.327--0.564} and mean Kendall
\(\tau\) ranges \textbf{0.653--0.895}. Removing native confidence
produces the largest stabilization, which reinforces rather than weakens
the calibration-fairness interpretation: comparison on unmatched native
confidence scales is a major source of apparent system selection.

Finally, the release includes
\texttt{external\_benchmarks/horizonbench\_adapter.py}, which consumes
the per-item JSONL result schema documented by the official HorizonBench
evaluation workflow. Standard HorizonBench output can be summarized by
overall/evolved/static accuracy; confidence-gated PersonaMetrica utility
is enabled only when an explicit \texttt{confidence} field is supplied.
The checked-in smoke test uses \textbf{synthetic schema fixtures only}.
No HorizonBench score, downloaded benchmark example, or frontier-model
result is claimed.

Full machine-readable outputs are in
\texttt{data/leaderboard\_seed\_uncertainty.json},
\texttt{data/protocol\_grid\_robustness.json}, and
\texttt{data/horizonbench\_adapter\_smoke.json}.

\section{NeurIPS 2026 paper
checklist}\label{neurips-2026-paper-checklist}

\begin{enumerate}
\def\labelenumi{\arabic{enumi}.}
\tightlist
\item
  \textbf{Claims --- Yes.} The Abstract, Introduction, Results, and
  Limitations explicitly scope claims to a controlled
  structured-evidence benchmark and separate executed results from
  future human/model experiments.
\item
  \textbf{Limitations --- Yes.} Section 10 enumerates
  structured-evidence, synthetic-user, binary-domain,
  designed-reliability, generator-coupling, external-model, human-study,
  utility, and multilingual limitations.
\item
  \textbf{Theory, assumptions, and proofs --- N/A.} The paper contains a
  probabilistic formulation but no theorem-level theoretical claim.
\item
  \textbf{Experimental result reproducibility --- Yes.} Code, fixed
  seeds, full generated data, exact experiment scripts, statistical
  scripts, figures, and reproduction logs are included in the anonymized
  bundle.
\item
  \textbf{Open access to data and code --- Not yet venue-compliant.} The
  anonymous supplementary bundle contains the code/data needed for the
  reported controlled results, but a persistent reviewer-accessible
  anonymous host has not been configured from this environment. This
  remains an explicit submission blocker and must be resolved before
  submission.
\item
  \textbf{Experimental setting/details --- Yes.} Sections 3--7 and
  Appendices B--D specify splits, hyperparameter selection, baselines,
  metric definitions, thresholds, costs, confidence calibration, and
  held-out seeds.
\item
  \textbf{Statistical significance --- Yes.} The main effects include
  bootstrap intervals over pre-specified simulator seeds and exact
  paired sign/sign-flip tests; the paper states what source of variation
  these quantify.
\item
  \textbf{Compute resources --- Yes.} Section 11 and
  \texttt{reports/COMPUTE\_RESOURCES.md} report CPU, RAM, wall time, and
  peak-memory examples.
\item
  \textbf{Code of Ethics --- Yes.} The executed benchmark uses synthetic
  users only; no deceptive human data collection or restricted personal
  data are used.
\item
  \textbf{Broader impacts --- Yes.} Section 11 discusses risks from
  persistent profiling, sensitive inference, and overconfident action,
  plus mitigations in the system artifact.
\item
  \textbf{Safeguards --- N/A for model release.} No high-risk pretrained
  model is released. Privacy, forgetting, permission, provenance, and
  verification hooks are included in the broader system artifact.
\item
  \textbf{Licenses --- Yes.} The project is released with a license and
  cites the research assets used for scientific comparison. No external
  dataset is repackaged into PersonaMetrica-Bench.
\item
  \textbf{New assets --- Yes.} The benchmark, generator, dataset card,
  Croissant metadata, benchmark specification, code, and limitations are
  documented alongside the release.
\item
  \textbf{Crowdsourcing/human subjects --- N/A for reported results.} No
  crowdsourced or human-subject data support the paper's current
  empirical claims. The unexecuted annotation protocol is included for
  future work only.
\item
  \textbf{IRB/equivalent approval --- N/A for reported results.} No
  human-subject experiment has been executed. Appropriate review would
  be obtained before running a study where required.
\end{enumerate}

\end{document}
